\section{Baselines and Diagnostic Models}

We compare three families of models.
\textbf{Null baselines} predict the majority label, primarily to expose the accuracy trap under skew.
\textbf{Semantic baselines} train TF-IDF or encoder representations with a linear classifier.
They test whether issue resolution can be recovered from generic text similarity.
\textbf{Structured diagnostic models} expose evidence slots and follow-up cues to a calibrator.
This model is deliberately simple: its purpose is to test whether revision-specific structure helps, not to maximize architecture novelty.

The issue-ledger and structured calibrator are designed as diagnostic interventions.
They encode cues such as concrete added experiments, softened claims, future-work deferrals, and response-only fixes.
Hard overrides are reported with an ablation so that the gain from explicit revision cues is visible.

\paragraph{Structured feature groups.}
The structured calibrator uses only observable fields from the issue sheet: request type cues in the review concern, response/revision cues for added evidence, broad-evaluation responses, theory or notation fixes, code-release promises, future-work deferrals, digit/count buckets, and the disagreement pattern among semantic baselines.
It also includes an issue-ledger label and rule as diagnostic features.
For ICLR 2024 clean-dev evaluation, we use leave-one-out training: each row is predicted by a classifier trained on all other labeled rows, with hard overrides disabled in the ablation row.
Transfer predictions train on the frozen labeled source sheet and are never used as labels for blind validation packets.
